\documentclass[11pt,a4paper]{report}

% XeLaTeX for Unicode support
\usepackage{fontspec}

% Use fonts with Kannada support
% Falls back gracefully if fonts not available
\setmainfont{Arial}
\newfontfamily\kannadafont{Noto Sans Kannada}[Script=Kannada]
\newcommand{\kannada}[1]{{\kannadafont #1}}

% Packages
\usepackage{geometry}
\geometry{margin=1in}
\usepackage{hyperref}
\usepackage{listings}
\usepackage{xcolor}
\usepackage{booktabs}
\usepackage{graphicx}
\usepackage{fancyhdr}
\usepackage{titlesec}

% Colors
\definecolor{codegreen}{rgb}{0,0.6,0}
\definecolor{codegray}{rgb}{0.5,0.5,0.5}
\definecolor{codepurple}{rgb}{0.58,0,0.82}
\definecolor{backcolour}{rgb}{0.97,0.97,0.95}
\definecolor{kapilacolor}{rgb}{0.72,0.53,0.04}

% Code listing style
\lstdefinestyle{mystyle}{
    backgroundcolor=\color{backcolour},
    commentstyle=\color{codegreen},
    keywordstyle=\color{codepurple},
    numberstyle=\tiny\color{codegray},
    stringstyle=\color{codegreen},
    basicstyle=\ttfamily\small,
    breakatwhitespace=false,
    breaklines=true,
    captionpos=b,
    keepspaces=true,
    numbers=left,
    numbersep=5pt,
    showspaces=false,
    showstringspaces=false,
    showtabs=false,
    tabsize=2,
    frame=single,
    framerule=0.5pt,
    rulecolor=\color{codegray}
}
\lstset{style=mystyle}

% Headers
\pagestyle{fancy}
\fancyhf{}
\fancyhead[L]{\leftmark}
\fancyhead[R]{\kannada{ಕಪಿಲ}}
\fancyfoot[C]{\thepage}

% Title formatting
\titleformat{\chapter}[display]
  {\normalfont\huge\bfseries\color{kapilacolor}}
  {\chaptertitlename\ \thechapter}{20pt}{\Huge}

% Document info
\title{
    \vspace{-1cm}
    {\Huge\bfseries\color{kapilacolor} \kannada{ಕಪಿಲ} (Kapila)}\\[0.5cm]
    {\Large A Kannada Programming Language}\\[1cm]
    {\large Language Manual \& Implementation Guide}
}
\author{
    \textbf{Architect and Author:} Prashanth Hebbar\\
    \texttt{hebbarp@gmail.com}\\[0.5cm]
    \textbf{Co-Authored By:} Claude Opus 4.5
}
\date{Version 0.6.0 \\ January 2026}

\begin{document}

\maketitle
\tableofcontents

% ============================================================
\chapter{Introduction}
% ============================================================

\section{What is Kapila?}

\kannada{ಕಪಿಲ} (Kapila) is a programming language that uses Kannada script natively. Unlike previous attempts that simply translated English keywords to Kannada, Kapila explores whether Kannada's grammatical structure can inspire new programming constructs.

Named after \textbf{Kapila} (\kannada{ಕಪಿಲ}), the ancient Vedic sage who founded the Samkhya school of philosophy—one of the oldest systems of Indian philosophical inquiry that emphasized systematic enumeration and logical analysis.

\section{Design Philosophy}

Kapila combines three programming paradigms:

\begin{itemize}
    \item \textbf{Forth's Soul}: Stack-based, concatenative core
    \item \textbf{Smalltalk's Face}: Clean, readable syntax
    \item \textbf{Perl's Tongue}: Do What I Mean (DWIM) flexibility
\end{itemize}

\section{Compilation Pipeline}

\begin{verbatim}
  Kannada Source Code
        |
        v
  +----------+     Tokens
  |  LEXER   | --> [WORD, NUMBER, ...]
  +----------+
        |
        v
  +----------+     Abstract Syntax Tree
  |  PARSER  | --> Program(statements=[...])
  +----------+
        |
        v
  +----------+     C Code
  | CODEGEN  | --> #include "kapila.h" ...
  +----------+
        |
        v
  +----------+     Native Executable
  | tcc/gcc  | --> program.exe
  +----------+
\end{verbatim}

% ============================================================
\chapter{Language Syntax}
% ============================================================

\section{Basic Elements}

\subsection{Numbers}

Kapila supports both Kannada and Arabic numerals:

\begin{lstlisting}
// Kannada numerals
\end{lstlisting}
\kannada{೧೨೩} \quad (123)\\
\kannada{೩.೧೪} \quad (3.14)

\begin{lstlisting}
// Arabic numerals also work
123
3.14
\end{lstlisting}

\subsection{Strings}

\begin{lstlisting}
"Hello, World!"
\end{lstlisting}
"\kannada{ನಮಸ್ಕಾರ ಪ್ರಪಂಚ}!"

\subsection{Booleans}

\begin{tabular}{ll}
\toprule
Kannada & English \\
\midrule
\kannada{ಸರಿ}, \kannada{ನಿಜ}, \kannada{ಹೌದು} & true \\
\kannada{ತಪ್ಪು}, \kannada{ಸುಳ್ಳು}, \kannada{ಇಲ್ಲ} & false \\
\bottomrule
\end{tabular}

\section{Arithmetic Operations}

\begin{tabular}{lll}
\toprule
Kannada & Symbol & Operation \\
\midrule
\kannada{ಕೂಡು} & + & Addition \\
\kannada{ಕಳೆ} & - & Subtraction \\
\kannada{ಗುಣಿಸು} & * & Multiplication \\
\kannada{ಭಾಗಿಸು} & / & Division \\
\kannada{ಶೇಷ} & \% & Modulo \\
\bottomrule
\end{tabular}

\vspace{0.5cm}
Example:
\begin{lstlisting}
// Postfix style (Forth-like)
5 3 + print.        // 8

// With Kannada words
\end{lstlisting}
\kannada{೫ ೩ ಕೂಡು ಮುದ್ರಿಸು.} \quad // 8

\section{Comparison Operations}

\begin{tabular}{lll}
\toprule
Kannada & Symbol & Operation \\
\midrule
\kannada{ಕಿರಿದು} & < & Less than \\
\kannada{ಹಿರಿದು} & > & Greater than \\
\kannada{ಸಮ} & = & Equal \\
\kannada{ಸಮನಲ್ಲ} & != & Not equal \\
\bottomrule
\end{tabular}

\section{Stack Operations}

\begin{tabular}{lll}
\toprule
Kannada & English & Effect \\
\midrule
\kannada{ನಕಲು} & dup & a -- a a \\
\kannada{ಬಿಡು} & drop & a -- \\
\kannada{ಅದಲುಬದಲು} & swap & a b -- b a \\
\kannada{ಮೇಲೆ} & over & a b -- a b a \\
\kannada{ತಿರುಗಿಸು} & rot & a b c -- b c a \\
\bottomrule
\end{tabular}

\section{Word Definitions}

Words (functions) are defined with the colon-double-danda syntax:

\begin{lstlisting}
// Define a word
square: dup * ||

// Use it
5 square print.    // 25
\end{lstlisting}

In Kannada:
\begin{verbatim}
\end{verbatim}
\kannada{ವರ್ಗ: ನಕಲು ಗುಣಿಸು ॥}\\
\kannada{೫ ವರ್ಗ ಮುದ್ರಿಸು.} \quad // 25

\section{Variables}

\begin{lstlisting}
x := 10.           // Assign
x print.           // Use
\end{lstlisting}

\section{Lists}

\begin{lstlisting}
[1 2 3]            // List literal
[1 2 3] length.    // 3
[1 2 3] first.     // 1
\end{lstlisting}

In Kannada:\\
\kannada{[೧ ೨ ೩] ಉದ್ದ ಮುದ್ರಿಸು.} \quad // 3

% ============================================================
\chapter{Code Generation}
% ============================================================

\section{Overview}

Kapila compiles to C code, which is then compiled to native executables using TinyCC or GCC.

\begin{verbatim}
Kapila Source --> C Code --> C Compiler --> Executable
   (.kpl)                     (tcc/gcc)      (.exe)
\end{verbatim}

\section{Why C Instead of LLVM?}

\begin{tabular}{lll}
\toprule
Aspect & C Backend & LLVM Backend \\
\midrule
Complexity & \textasciitilde300 lines & \textasciitilde1000+ lines \\
Dependencies & Just tcc/gcc & llvmlite + LLVM (\textasciitilde500MB) \\
Debugging & Read the .c file & Binary IR \\
Distribution & 18MB & 500MB+ \\
\bottomrule
\end{tabular}

\section{Generated Code Structure}

For the Kapila program:
\begin{verbatim}
\end{verbatim}
\kannada{ವರ್ಗ: ನಕಲು ಗುಣಿಸು ॥}\\
\kannada{೫ ವರ್ಗ ಮುದ್ರಿಸು.}

\vspace{0.3cm}
The generator produces:

\begin{lstlisting}[language=C]
/* Generated by Kapila Compiler */
#include "kapila.h"

void word_varga(void) {
    dup_op();
    mul_op();
}

int main() {
    kapila_init();

    push_int(5);
    word_varga();
    println_op();

    kapila_cleanup();
    return 0;
}
\end{lstlisting}

\section{Runtime Library}

The runtime library (\texttt{runtime/kapila.h} and \texttt{runtime/kapila.c}) provides:

\begin{itemize}
    \item Value type with union (int, float, bool, string, list)
    \item Stack operations (push, pop, peek)
    \item Arithmetic operations
    \item Comparison operations (with float support)
    \item String operations (UTF-8 aware length)
    \item List operations
    \item Memory management (arena allocator)
\end{itemize}

% ============================================================
\chapter{The Kapila Compiler}
% ============================================================

\section{Usage}

\begin{lstlisting}[language=bash]
# Show version
kapilac -v

# Compile and run
kapilac program.kpl -r

# Output C code
kapilac program.kpl -o output.c

# Compile from string
kapilac -c '5 3 + print.'
\end{lstlisting}

\section{Command-Line Options}

\begin{tabular}{ll}
\toprule
Option & Description \\
\midrule
\texttt{input.kpl} & Input source file \\
\texttt{-o FILE} & Output C file \\
\texttt{-r, --run} & Compile and run \\
\texttt{-c CODE} & Compile from string \\
\texttt{-k, --keep} & Keep temporary files \\
\texttt{-v, --version} & Show version \\
\bottomrule
\end{tabular}

\section{Bundled TinyCC}

The standalone distribution includes TinyCC (Tiny C Compiler):

\begin{itemize}
    \item Size: \textasciitilde1MB (vs GCC's 500MB+)
    \item Fast compilation
    \item No external dependencies
\end{itemize}

\section{Distribution Structure}

\begin{verbatim}
kapila/
  kapilac.exe     # Compiler
  kapila.exe      # REPL/Interpreter
  runtime/        # C runtime library
    kapila.h
    kapila.c
  tcc/            # Bundled TinyCC
  examples/       # Sample programs
\end{verbatim}

% ============================================================
\chapter{Future: LLVM Backend}
% ============================================================

\textbf{Status: PLANNED -- Not yet implemented}

\section{Motivation}

A future version may include an LLVM backend for:

\begin{itemize}
    \item Advanced optimizations (-O3, LTO)
    \item JIT compilation for REPL
    \item WebAssembly output
    \item Full debug information
\end{itemize}

\section{Proposed Architecture}

\begin{verbatim}
                kapilac
                   |
            +------+------+
            |             |
            v             v
      +---------+   +-----------+
      | C Back- |   | LLVM Back-|
      | end     |   | end       |
      +---------+   +-----------+
            |             |
            v             v
         tcc/gcc        LLVM
            |             |
            v             v
         .exe         .exe/.wasm
\end{verbatim}

\section{LLVM IR Example}

For \kannada{೫ ೩ ಕೂಡು ಮುದ್ರಿಸು.}, the LLVM IR would be:

\begin{lstlisting}
define i32 @main() {
entry:
    call void @push_int(i64 5)
    call void @push_int(i64 3)
    call void @add_op()
    call void @println_op()
    call void @kapila_cleanup()
    ret i32 0
}
\end{lstlisting}

% ============================================================
\chapter{Quick Reference}
% ============================================================

\section{Syntax Cheatsheet}

\begin{lstlisting}
// Statement endings
statement.           // End with period
word: body ||        // Word definition ends with ||

// Assignment
x := 10.

// Arithmetic (postfix)
5 3 +               // 8
\end{lstlisting}

\section{Kannada Vocabulary}

\subsection{Arithmetic}
\begin{tabular}{ll}
\kannada{ಕೂಡು} & add \\
\kannada{ಕಳೆ} & subtract \\
\kannada{ಗುಣಿಸು} & multiply \\
\kannada{ಭಾಗಿಸು} & divide \\
\end{tabular}

\subsection{Comparison}
\begin{tabular}{ll}
\kannada{ಕಿರಿದು} & less than \\
\kannada{ಹಿರಿದು} & greater than \\
\kannada{ಸಮ} & equal \\
\end{tabular}

\subsection{Stack}
\begin{tabular}{ll}
\kannada{ನಕಲು} & dup \\
\kannada{ಬಿಡು} & drop \\
\kannada{ಅದಲುಬದಲು} & swap \\
\end{tabular}

\subsection{I/O}
\begin{tabular}{ll}
\kannada{ಮುದ್ರಿಸು} & print \\
\kannada{ಉದ್ದ} & length \\
\kannada{ಮೊದಲ} & first \\
\end{tabular}

% ============================================================
\chapter{Example Programs}
% ============================================================

\section{Hello World}

\begin{verbatim}
"Hello, World!" print.
\end{verbatim}

In Kannada:\\
\kannada{"ನಮಸ್ಕಾರ ಪ್ರಪಂಚ!" ಮುದ್ರಿಸು.}

\section{Square Function}

\begin{verbatim}
square: dup * ||
5 square print.    // 25
\end{verbatim}

In Kannada:\\
\kannada{ವರ್ಗ: ನಕಲು ಗುಣಿಸು ॥}\\
\kannada{೫ ವರ್ಗ ಮುದ್ರಿಸು.}

\section{Arithmetic Report}

\begin{verbatim}
// Header
"=== Calculation ===" print.

// Operations
12 8 + print.      // 20
7 6 * print.       // 42

// Using words
double: 2 * ||
5 double print.    // 10

"Done!" print.
\end{verbatim}

% ============================================================
\appendix
\chapter{Installation}
% ============================================================

\section{Standalone Distribution}

\begin{enumerate}
    \item Download \texttt{kapila-windows-x64.zip} from GitHub Releases
    \item Extract to any folder
    \item Run \texttt{kapilac.exe} or \texttt{kapila.exe}
\end{enumerate}

No Python or C compiler installation required.

\section{From Source}

\begin{lstlisting}[language=bash]
git clone https://github.com/hebbarp/kapila
cd kapila
python kapila.py           # REPL
python kapilac.py file.kpl -r  # Compile & run
\end{lstlisting}

Requirements: Python 3.10+, GCC or TinyCC

% ============================================================
\chapter{License and Credits}
% ============================================================

\section{Author}

\textbf{Architect and Author:} Prashanth Hebbar (\texttt{hebbarp@gmail.com})

\textbf{Co-Authored By:} Claude Opus 4.5

\section{Repository}

\url{https://github.com/hebbarp/kapila}

\section{Acknowledgments}

\begin{itemize}
    \item Named after Kapila, the ancient Vedic sage
    \item Unicode handling adapted from Pampa project
    \item TinyCC by Fabrice Bellard
\end{itemize}

\end{document}
